\documentclass[11pt,a4paper]{article}
\usepackage[margin=2.5cm]{geometry}
\usepackage{amsmath,amssymb}
\usepackage{booktabs}
\usepackage{enumitem}
\usepackage{xcolor}
\usepackage[hidelinks]{hyperref}
\setlist{itemsep=2pt,topsep=4pt}
\newcommand{\code}[1]{\texttt{#1}}
\newcommand{\meas}{\textcolor{teal}{\textsf{[measured]}}}
\newcommand{\hyp}{\textcolor{orange!80!black}{\textsf{[hypothesis]}}}
\newcommand{\E}{\mathbb{E}}
\newcommand{\Gt}{\widetilde{G}}

\title{When is the Gram matrix of MGD unstable, and what we do about it\\[2pt]
\large A step-by-step explanation, with the measurements of the 30 September run}
\author{}
\date{1 October 2026}

\begin{document}
\maketitle

\section*{What this note is for}
At every time step, MGD moves the walkers by solving a linear system whose matrix is the
\emph{Gram matrix} $G$ of the gradients of the statistics. If $G$ is ``unstable''
(the word used in numerics is \emph{ill-conditioned}), the solution can become huge or
meaningless. This note explains, without assuming any background in numerical linear
algebra:
\begin{enumerate}
  \item where $G$ comes from (Sections~\ref{sec:mgd}--\ref{sec:where});
  \item the two ways in which $G$ can make the solve fail (Section~\ref{sec:two});
  \item the three remedies used in the code: rescaling, ridge, mask (Section~\ref{sec:remedies});
  \item what we measured along the transport in the run of 30 September (Section~\ref{sec:measured});
  \item what this does and does not explain (Section~\ref{sec:scope}).
\end{enumerate}
\meas\ marks a number computed on the run or in a controlled test; \hyp\ marks an
interpretation that has not been tested directly.

%=====================================================================================
\section{MGD in one page}\label{sec:mgd}

\paragraph{Statistics.} A signal is a vector $x\in\mathbb R^M$ (here $M=256$). We describe it by
$r$ real-valued \emph{statistics} $\phi_1(x),\dots,\phi_r(x)$ (here $r=281$): $L^6$ norms of
wavelet coefficients, region statistics, scattering coefficients, \dots. We write
$\phi=(\phi_1,\dots,\phi_r)$. Each $\phi_n$ has a gradient $\nabla\phi_n(x)\in\mathbb R^M$: the
direction in which one should move the signal $x$ to increase $\phi_n$ fastest.

\paragraph{Walkers and their moments.} MGD evolves $N$ signals $x^1,\dots,x^N$, the
\emph{walkers} (here $N=8500$), from time $t=0$ to $t=1$ on a grid $t_0<t_1<\dots<t_{n_t}$.
The walkers at step $k$ are written $x_k^i$. Their \emph{empirical moments} are the averages
\[
  \bar\phi_n(x_k)=\frac1N\sum_{i=1}^N\phi_n(x_k^i),\qquad n=1,\dots,r .
\]

\paragraph{Target moments.} Let $x_1^1,\dots,x_1^P$ be the data and $x_0^1,\dots,x_0^P$ independent
Gaussian noise. The \emph{interpolant} between noise and data is
\[
  I_t^j=\cos\!\Big(\frac{\pi t}{2}\Big)x_0^j+\sin\!\Big(\frac{\pi t}{2}\Big)x_1^j,\qquad j=1,\dots,P,
\]
pure noise at $t=0$, the data at $t=1$. The \emph{target moments} at time $t$ are
\[
  m^\star_n(t)=\frac1P\sum_{j=1}^P\phi_n\big(I_t^j\big).
\]
No model of the data is needed: the targets are computed from data and noise.
The aim of MGD is that, at every time, the walkers have the target moments,
$\bar\phi_n(x_k)\approx m^\star_n(t_k)$ for all $n$, while being otherwise as random
(maximum entropy) as possible.

\paragraph{One time step.} With $h=t_{k+1}-t_k$, each step has two parts (code:
\code{iteration\_step\_projection}).
\begin{itemize}
  \item \textbf{Predictor.} Move the walkers along the gradients so that their moments follow
    the motion of the target, and add noise of strength $\sigma$ (here $\sigma=3.5$):
    \[
      y_k^i=x_k^i+h\sum_{m=1}^r\eta_m\nabla\phi_m(x_k^i)+\sqrt{2h}\,\sigma\,\xi^i,
      \qquad \xi^i\sim\mathcal N(0,I).
    \]
  \item \textbf{Corrector.} Move the walkers once more along the gradients to remove the
    remaining gap to the target at $t_{k+1}$:
    \begin{equation}
      x_{k+1}^i=y_k^i+\sum_{m=1}^r\theta_m\nabla\phi_m(y_k^i).
      \label{eq:corrector}
    \end{equation}
\end{itemize}
The unknowns are the coefficients $\eta=(\eta_1,\dots,\eta_r)$ and $\theta=(\theta_1,\dots,\theta_r)$,
one per statistic. (The $\theta$ saved to disk is divided by $h\sigma^2$; this does not
matter here.) Both are obtained by solving a linear system with the Gram matrix. We now see why.

%=====================================================================================
\section{Where the Gram matrix comes from}\label{sec:where}

Take the corrector \eqref{eq:corrector}. How much does it change the moment $\bar\phi_n$?
For a small move, a function changes by (gradient)$\cdot$(displacement); this is the
first-order Taylor expansion:
\[
  \phi_n(x_{k+1}^i)\;\approx\;\phi_n(y_k^i)+\nabla\phi_n(y_k^i)\cdot\big(x_{k+1}^i-y_k^i\big)
  =\phi_n(y_k^i)+\sum_{m=1}^r\theta_m\,\nabla\phi_n(y_k^i)\cdot\nabla\phi_m(y_k^i).
\]
Averaging over the walkers,
\[
  \bar\phi_n(x_{k+1})-\bar\phi_n(y_k)\;\approx\;\sum_{m=1}^r G_{nm}\,\theta_m,
  \qquad
  \boxed{\,G_{nm}=\frac1N\sum_{i=1}^N\nabla\phi_n(y_k^i)\cdot\nabla\phi_m(y_k^i)\,}
\]
$G$ is the \textbf{Gram matrix of the gradients}: entry $(n,m)$ is the inner product of the
gradients of statistics $n$ and $m$, averaged over the walkers. We want the change to close
the gap to the target, $b_n=m^\star_n(t_{k+1})-\bar\phi_n(y_k)$. This gives one linear equation
per statistic,
\begin{equation}
  G\,\theta=b .
  \label{eq:system}
\end{equation}
The predictor gives the same kind of system, $G\,\eta=\dot m^\star(t_k)$, with $G$ evaluated on
$x_k$ and the right-hand side equal to the speed at which the target moments move.

\paragraph{Reading $G$.}
\begin{itemize}
  \item The \textbf{diagonal entry} $G_{nn}=\frac1N\sum_i|\nabla\phi_n(y_k^i)|^2$ says how strongly
    statistic $n$ reacts when the walkers move. If $\nabla\phi_n=0$ on every walker, no move of
    the form \eqref{eq:corrector} can change $\bar\phi_n$.
  \item The \textbf{off-diagonal entry} $G_{nm}$ says how much moving along $\nabla\phi_m$ also changes
    $\phi_n$. If the two gradients are orthogonal on average over the walkers, $G_{nm}=0$: the two
    statistics can be adjusted independently. If they are almost parallel, adjusting one also
    adjusts the other.
\end{itemize}

%=====================================================================================
\section{Two ways in which the solve can fail}\label{sec:two}

Solving \eqref{eq:system} means finding the coefficients $\theta$ that produce the requested
changes $b$. It fails when some requested change can only be produced by an enormous move.
This happens in two distinct ways.

\subsection{Failure 1: a statistic that the walkers cannot move}\label{sec:fail1}
Take the simplest case: the gradient of statistic $n$ is orthogonal, on average over the walkers,
to the gradients of all the others ($G_{nm}=0$ for every $m\neq n$). Row $n$ of
\eqref{eq:system} then reads $G_{nn}\theta_n=b_n$, so
\[
  \theta_n=\frac{b_n}{G_{nn}} .
\]
The size of the resulting move is $\frac1N\sum_i|\theta_n\nabla\phi_n(y_k^i)|^2=\theta_n^2G_{nn}=b_n^2/G_{nn}$.
If $G_{nn}\to0$ and the gap $b_n$ does not shrink at least as fast as $\sqrt{G_{nn}}$, the move
diverges: we ask for a change of a moment that the walkers can barely influence. Nothing in the
method forces $b_n$ to shrink that fast: $b_n$ is the difference of two averages over finitely many
samples, and their sampling noise does not depend on $G_{nn}$.

\textbf{Where it happens.} Region statistics (\code{Scalar\_*\_gaussianK}) measure how much of a
wavelet channel falls in a given range of amplitudes. Early in the transport the walkers are close to
Gaussian noise, and some amplitude ranges contain almost no walker. There $\nabla\phi_n\approx0$
and $G_{nn}\approx0$ (on the walkers, $G_{nn}$ reached $10^{-16}$--$10^{-8}$ times its data value
\meas). The walker blow-up of job 268295 at step 163 (28 September) was diagnosed as this, and the
rerun with the mask of Section~\ref{sec:remedies} completed \meas.

\subsection{Failure 2: two statistics that move the walkers in the same way}\label{sec:fail2}
Now take two statistics $a$ and $b$ whose gradients are almost the same (after normalising both to
the same size). With unit diagonal, their $2\times2$ block of $G$ is
\[
  G=\begin{pmatrix}1&1-\delta\\ 1-\delta&1\end{pmatrix},\qquad \delta \text{ small}.
\]
Solving $G\theta=b$ by hand (add and subtract the two equations):
\[
  \theta_a+\theta_b=\frac{b_a+b_b}{2-\delta},\qquad
  \theta_a-\theta_b=\frac{b_a-b_b}{\delta}.
\]
The \emph{sum} $\theta_a+\theta_b$ is well determined. The \emph{difference} is divided by
$\delta$. Both statistics push the walkers in (almost) the same direction, so the walkers only
``see'' their sum, and the difference is decided by a tiny $\delta$. Whatever is in $b_a-b_b$,
including sampling noise (targets and walker moments are averages over finitely many samples) and
rounding errors, is multiplied by $1/\delta$. For instance, with $\delta=10^{-7}$, a difference
$b_a-b_b=10^{-6}$ gives $\theta_a-\theta_b=10$. Whether the result is actually large depends on
how small $b_a-b_b$ happens to be; the point is that the solve has no control over it.

Note that \emph{neither statistic is weak}: both have $G_{aa}=G_{bb}=1$. Failure~2 is invisible on
the diagonal; it is in the off-diagonal entry $1-\delta$.

\subsection{The general picture: eigen-directions and the condition number}
Both failures are special cases of one picture. $G$ is symmetric, so it has $r$ orthogonal
\emph{eigen-directions} $u_1,\dots,u_r$ (each is a combination of the $r$ coefficients) with
eigenvalues $\mu_1\le\mu_2\le\dots\le\mu_r$, all $\ge0$. The solution of \eqref{eq:system} is
\[
  \theta=\sum_{i=1}^r\frac{u_i\cdot b}{\mu_i}\,u_i .
\]
Read it as: the eigenvalue $\mu_i$ says how much the moments change when the coefficients move
along $u_i$; to obtain the requested change, the solve divides by $\mu_i$. A \textbf{small eigenvalue}
is a combination of coefficients that the walkers barely see, and whatever part of $b$ lies along it
is blown up. In the $2\times2$ example, the eigenvalues are $2-\delta$ (direction ``$a+b$'')
and $\delta$ (direction ``$a-b$''). In Failure~1, the small eigenvalue is $G_{nn}$ itself, with
direction ``statistic $n$ alone''.

The \textbf{condition number} is the ratio of the largest to the smallest eigenvalue,
\[
  \kappa=\frac{\mu_{\max}}{\mu_{\min}} .
\]
Rule of thumb: errors in $G$ or $b$ can be amplified by up to a factor $\kappa$ in the solution,
i.e.\ one can lose about $\log_{10}\kappa$ significant digits. A computer number in single precision
(\code{float32}) carries about 7 digits, in double precision (\code{float64}) about 16. So when
$\kappa$ approaches $10^7$, a matrix stored in \code{float32} no longer determines the solution
along its weakest directions.

%=====================================================================================
\section{The three remedies in the code}\label{sec:remedies}

\paragraph{1. Rescaling (Jacobi).} The statistics have different units, so the raw $G$ mixes large and
small numbers for no physical reason. The code first divides each statistic by $\sqrt{G_{nn}}$:
\[
  \Gt_{nm}=\frac{G_{nm}}{\sqrt{G_{nn}G_{mm}}},\qquad \Gt_{nn}=1 .
\]
$\Gt_{nm}$ is then the \emph{cosine} between the two gradients, each seen as one long vector over all
walkers and all positions: $0$ if they are orthogonal on average, close to $\pm1$ if almost parallel. All condition numbers below are for $\Gt$. Rescaling
removes differences of units but cannot fix either failure: in Failure~2 the cosine stays
$1-\delta$, and in Failure~1 the division is by the vanishing $\sqrt{G_{nn}}$ itself.

\paragraph{2. Ridge $\varepsilon$.} The code solves $(\Gt+\varepsilon I)\tilde\theta=\tilde b$ with
$\varepsilon=10^{-4}$ (\code{--regularization}). In the eigen-picture,
\[
  \tilde\theta=\sum_i\frac{u_i\cdot\tilde b}{\mu_i+\varepsilon}\,u_i :
\]
directions with $\mu_i\gg\varepsilon$ are untouched; directions with $\mu_i\ll\varepsilon$ get the
coefficient $(u_i\cdot\tilde b)/\varepsilon$ instead of $(u_i\cdot\tilde b)/\mu_i$. So the ridge
\emph{caps} Failure~2, at the price that the coefficient along these directions is set by the
arbitrary $\varepsilon$ and not by the data.

The ridge does \emph{not} cure Failure~1. Undoing the rescaling, the solved system is
$(G+\varepsilon\,\mathrm{diag}\,G)\,\theta=b$. In the uncoupled case this gives
$\theta_n=b_n/\big((1+\varepsilon)G_{nn}\big)$: the ridge is proportional to $G_{nn}$ and vanishes with it.

\paragraph{3. The mask (for Failure 1).} At each step we keep only the statistics that can still act on
the walkers,
\[
  L=\big\{\,n:\ G_{nn}>\tau_n\,\big\},
\]
with $\tau_n=10^{-6}$ for the region statistics and $\tau_n=0$ for all others. Each statistic is
normalised so that $\E_{\rm data}|\nabla\phi_n|^2=1$, so $\tau_n=10^{-6}$ means one millionth of its
value on the data. In practice, row and column $n$ of $G$ are removed for every $n\notin L$: we solve
the smaller system $(G_{LL}+\varepsilon\,\mathrm{diag}\,G_{LL})\,\theta_L=b_L$ and set $\theta_n=0$
for the others. This barely affects the live statistics, because by the Cauchy--Schwarz inequality
$|G_{nm}|\le\sqrt{G_{nn}G_{mm}}$: when $G_{nn}$ is close to zero, the whole row and column $n$ are
close to zero as well. Setting $\theta_n=0$ removes almost no control, since the walkers could barely
move $\bar\phi_n$ anyway; the price is that $\bar\phi_n$ is not corrected while the statistic is masked. When the amplitude range fills up again, $G_{nn}>\tau_n$ and the statistic
re-enters $L$ automatically.

The mask looks only at the diagonal, so it cannot see Failure~2: two near-copies each have a normal
diagonal.

%=====================================================================================
\section{What we measured along the transport (run of 30 September)}\label{sec:measured}

Job 382722: turbulence, seed 900, two-phase time grid ($n_t=60\,000$), $N=8500$, $r=281$,
$\varepsilon=10^{-4}$. The run logged the eigenvalues of $\Gt$ every 10 steps (6000 logged steps),
for both solves; the figures are in \code{turbulence/gram\_along\_transport.ipynb}. The run
finished normally \meas.

\subsection{Failure 1: masking}
Only 5 of the 281 statistics were ever masked, all region statistics, all before $t=0.15$ \meas:
\begin{center}\small
\begin{tabular}{lcc}
\toprule
statistic & masked during & masked steps\\
\midrule
\code{Scalar\_morlet\_gaussianK[21]} & $t<0.021$ & 236\\
\code{Scalar\_psi\_gaussianK[55]} & $t<0.044$ (short breaks at $t\approx0.021$ and $0.0435$) & 478\\
\code{Scalar\_psi\_gaussianK[33]} & $t<0.055$ (short break at $t\approx0.021$) & 603\\
\code{Scalar\_morlet\_gaussianK[30]} & $t<0.148$ & 1646\\
\code{Scalar\_psi\_gaussianK[77]} & $t<0.149$ & 1657\\
\bottomrule
\end{tabular}
\end{center}
Seven statistics have a coefficient jump, defined as $|\theta_n|$ above 100 times its own running
median. For all seven, the largest jump is at $t\approx0.0213$--$0.0215$ \meas. This is the moment
when \code{Scalar\_morlet\_gaussianK[21]} is switched back on, with little margin above its threshold,
and \code{Scalar\_psi\_gaussianK[33]} and \code{[55]} are briefly switched on as well \meas. That the
re-activation causes the jumps is likely but not proven \hyp. Four of the seven statistics have a
second, smaller jump episode, whose time I have not located.

\subsection{Failure 2: near-copies}
The four smallest eigenvalues of $\Gt$ (corrector), median over each time window \meas:
\begin{center}\small
\begin{tabular}{lccccc}
\toprule
$t$ window & $\mu_1$ & $\mu_2$ & $\mu_3$ & $\mu_4$ & $\kappa=\mu_{\max}/\mu_1$\\
\midrule
$[0,\,0.1)$     & $6.0\cdot10^{-7}$ & $4.4\cdot10^{-6}$ & $2.6\cdot10^{-5}$ & $2.7\cdot10^{-4}$ & $2.5\cdot10^{7}$\\
$[0.1,\,0.5)$   & $1.9\cdot10^{-7}$ & $5.0\cdot10^{-7}$ & $5.3\cdot10^{-6}$ & $3.5\cdot10^{-4}$ & $6.2\cdot10^{7}$\\
$[0.5,\,0.9)$   & $1.6\cdot10^{-7}$ & $6.7\cdot10^{-7}$ & $1.2\cdot10^{-5}$ & $5.1\cdot10^{-4}$ & $7.2\cdot10^{7}$\\
$[0.99,\,0.999)$& $1.5\cdot10^{-7}$ & $6.7\cdot10^{-7}$ & $1.3\cdot10^{-5}$ & $5.1\cdot10^{-4}$ & $1.3\cdot10^{8}$\\
$[0.999,\,1)$   & $1.4\cdot10^{-7}$ & $6.7\cdot10^{-7}$ & $1.3\cdot10^{-5}$ & $5.0\cdot10^{-4}$ & $1.5\cdot10^{8}$\\
\bottomrule
\end{tabular}
\end{center}
Three eigenvalues are below the ridge $\varepsilon=10^{-4}$ at almost every logged step (two or four
at a few steps near $t=0$), and $\mu_4$ is a factor 10--70 above $\mu_3$ \meas. So $\Gt$ has
\textbf{three directions far weaker than all the others}.

\paragraph{Which statistics.} For the smallest direction, the log records the statistics with the
largest weight. At 96\% of the steps it is the pair $L_6[7]$ / $L_6^\psi[20]$; at 3\% the pair
$L_6[8]$ / $L_6^\psi[21]$; at a few steps the pair \code{Scalar\_psi\_gaussianK[77]} /
\code{Scalar\_morlet\_gaussianK[30]} \meas. A search for numerically null directions on the saved
systems (25 times checked) finds, at $t=0.9985$, the direction $0.71\,[L_6[7]]-0.71\,[L_6^\psi[20]]$:
exactly the ``$a-b$'' direction of Section~\ref{sec:fail2} \meas. The log records the composition of the
smallest direction only. That $\mu_1$ and $\mu_2$ are the two $L_6$ pairs is therefore an inference from
which pairs show up as smallest, and the third direction ($\mu_3$) is probably the Scalar pair; neither
has been checked directly \hyp.

\paragraph{Why these pairs exist.} $L_6$ is computed on a wavelet bank with $Q=1$, $L_6^\psi$ on a
bank with $Q=3$. In the usual construction of such banks, $Q$ sets the bandwidth \emph{relative} to the
centre frequency. But the grid has $M=256$ points, so frequencies only come in integer bins. At the
coarsest scales, both filters of each pair put essentially all their weight on the same one or two bins
(bin~1 for the first pair, bin~0 for the second) \meas, so they are almost identical. Measured as
$1-|\cos|$ between the filters, they differ by $5.8\cdot10^{-4}$ and $1.5\cdot10^{-3}$; between the
gradients of the corresponding $L_6$ statistics, on synthetic signals and in \code{float64}, the
difference is only $4\cdot10^{-7}$--$1.4\cdot10^{-5}$ \meas. The code already removes duplicate filters
\emph{inside} the $Q=3$ bank, but never compares the two banks with each other.

\paragraph{How reliable are these numbers.} $G$ is accumulated in \code{float32} (about 7 digits) and
only then converted to \code{float64} for the solve. Rounding each entry to 7 digits moves the eigenvalues
by an amount that, for a $281\times281$ matrix, can reach roughly $10^{-7}$--$10^{-6}$ \hyp\ (order-of-magnitude
estimate, not measured). At 5\% of the logged steps the smallest eigenvalue indeed comes out slightly
\emph{negative} ($\approx-10^{-7}$), which is impossible for a Gram matrix \meas. So:
\begin{itemize}
  \item $\mu_1$ and $\mu_2$ ($\sim10^{-7}$--$10^{-6}$) are at the precision limit. Their order of magnitude
    is consistent with the \code{float64} test on synthetic signals above, but their values, and hence
    $\kappa\sim10^7$--$10^8$, cannot be trusted precisely. What is certain is that these directions are
    far below the ridge.
  \item $\mu_3$ and $\mu_4$ ($\sim10^{-5}$ and $\sim5\cdot10^{-4}$) are well above the precision limit.
\end{itemize}
Only building $G$ in \code{float64} would give the true $\mu_1$, $\mu_2$.

\subsection{Without the near-copies}
Removing $k$ statistics from $G$ (deleting their rows and columns) cannot make the new smallest
eigenvalue larger than the old $\mu_{k+1}$; this is a standard result, \emph{eigenvalue interlacing}.
The bound is approximately reached when the removed statistics are the ones that make up the weak
directions. The largest eigenvalue barely changes when 2--3 of 281 statistics are removed. So the table
above gives, approximately, the best condition number that removing statistics can achieve:
\begin{center}\small
\begin{tabular}{lcc}
\toprule
statistics removed & new smallest eigenvalue & best $\kappa$\\
\midrule
none & $\mu_1\approx1.5\cdot10^{-7}$ & $\sim10^{8}$ (at the precision limit)\\
2 statistics (one per $L_6$ pair)$^*$ & $\mu_3\approx1.2\cdot10^{-5}$ & $\sim10^{6}$\\
3 statistics (also the third direction) & $\mu_4\approx5\cdot10^{-4}$ & $\sim2$--$5\cdot10^{4}$\\
\bottomrule
\end{tabular}\\[2pt]
{\footnotesize $^*$ assuming $\mu_1,\mu_2$ are the two $L_6$ pairs (see above).}
\end{center}
The last value stays between $2\cdot10^4$ and $5\cdot10^4$ in every time window, from $t=0$ to $t=1$
\meas. Near $t=1$ the walkers should be distributed approximately like the data, so this suggests that
$\sim3\cdot10^4$ is also the conditioning of the full set of 281 statistics on the data, i.e.\ a property
of the statistics rather than something the transport creates \hyp. A direct check is to compute $\Gt$ on
the data. At that level the ridge still acts, but mildly: along the weakest remaining direction the
solution is multiplied by $\mu_4/(\mu_4+\varepsilon)\approx0.7$--$0.8$, and by almost $1$ along all the
others. \code{float32} has a comfortable margin.

\paragraph{Correction to the note of 30 September.} The value $\kappa\approx12$ quoted there was for the
73 \code{Scalar\_psi\_gaussianK} statistics alone, on the data. The full 281-statistic Gram contains
near-dependencies \emph{between} families (e.g.\ $L_6$ on $Q=1$ and on $Q=3$ filters) that a single family
cannot show.

%=====================================================================================
\section{What this explains, and what it does not}\label{sec:scope}

\begin{itemize}
  \item \textbf{Explained.} The crash of 28 September (Failure~1, cured by the mask), and why $\Gt$ has a
    very large condition number (Failure~2: three weak directions, whose coefficients are set by
    $\varepsilon$; the weakest certainly, and most likely the second, come from near-copy filters).
  \item \textbf{Not explained.} The moments that are still not matched at the end of the run: 11
    moments miss their target by more than 1\% at the last time, all of them region statistics, with errors
    that first exceed 1\% between $t=0.957$ and $t=0.995$ \meas. None of these statistics is ever masked, and
    none is ever among the main components of the smallest eigen-direction \meas. We have no evidence
    linking them to the conditioning of $G$; but the composition of the second and third weakest
    directions is unknown, so this is not excluded. Their cause is under investigation.
\end{itemize}

\paragraph{About the supervisor's suggestion (conjugate gradient with early stopping).} Stopping when
$\|G\theta-b\|$ falls below the statistical noise of $b$ is another way of regularising, in which the noise,
not a hand-picked $\varepsilon$, sets the threshold (the \emph{discrepancy principle}). Replayed on the
logged systems, at 5999 of 6000 steps the gap $b$ is already below its estimated noise level before any
correction ($\theta=0$), so the method would stop immediately and not correct at all \meas. This holds with both the plain and the
noise-weighted norm. The noise estimate itself may be too large, for instance if the noise of the target is
not independent from step to step \hyp; until that is checked, the criterion cannot be used as is.

\section*{Proposed next steps}
\begin{enumerate}
  \item Remove the cross-bank near-copies when the filters are built: drop the $Q=1$ filters 7 and 8
    (difference to a $Q=3$ filter below $10^{-2}$; all other $Q=1$ filters are at least $7\cdot10^{-2}$ away).
  \item Identify the third near-singular direction, from the eigenvector of $\mu_3$ on the saved systems.
  \item Optionally, build $G$ in \code{float64}, to measure the smallest eigenvalues beyond the
    \code{float32} limit.
\end{enumerate}

\end{document}
